% 李群（综述）
% license CCBYSA3
% type Wiki

本文根据 CC-BY-SA 协议转载翻译自维基百科\href{https://en.wikipedia.org/wiki/Lie_group}{相关文章}

在数学中，李群（Lie group，发音为 /liː/，读作“LEE”）是既是群又是可微流形的结构，其特点是群的乘法运算和取逆运算都是可微的。

流形是局部上类似于欧几里得空间的空间；而群则定义了一种抽象的二元运算及其必须具备的附加性质，使其能被看作一种“抽象变换”。例如，群可以通过乘法与逆元运算（对应于除法），或等价地，通过加法与减法来体现。将这两种概念结合，就得到了连续群，其中点的乘法和取逆运算是连续的。如果乘法和取逆进一步是光滑（可微）的，那么这样的连续群就是李群。

李群为连续对称性的概念提供了一种自然模型，其中著名的例子就是圆群。旋转一个圆就是连续对称性的实例。对于圆的任意一次旋转，都存在同样的对称性；而这些旋转的组合（拼接）构成了圆群，这是李群的典型例子。李群在现代数学和物理的许多领域都有广泛的应用。

李群最初是在研究矩阵子群
$G \subseteq \mathrm{GL}_n(\mathbb{R})$ 或 $\mathrm{GL}_n(\mathbb{C})$ 时被发现的，这里的 $\mathrm{GL}_n(\mathbb{R})$ 与 $\mathrm{GL}_n(\mathbb{C})$ 分别表示实数域 $\mathbb{R}$ 或复数域 $\mathbb{C}$ 上的 $n \times n$ 可逆矩阵群。如今，这些群被称为**经典群**，因为李群的概念已经远远超出了这些最初的范围。

李群的名称源自挪威数学家索弗斯·李（Sophus Lie，1842–1899），他奠定了连续变换群理论的基础。李引入李群的最初动机，是为了刻画微分方程的连续对称性，这与伽罗瓦理论中利用有限群来刻画代数方程的离散对称性有着类似的思想。
\subsection{历史}
索弗斯·李认为1873–1874 年冬天是他连续群理论的诞生时期。[2] 然而，托马斯·霍金斯则认为，应当追溯到“李在 1869 年秋天至 1873 年秋天这四年间的惊人研究活动”，这才真正促成了该理论的建立。[2]李的一些早期思想是在与费利克斯·克莱因的密切合作中发展起来的。从 1869 年 10 月到 1872 年，李几乎每天都与克莱因见面：从 1869 年 10 月底到 1870 年 2 月底在柏林，其后两年则在巴黎、哥廷根和埃尔兰根。[3]李本人声称，到 1884 年时，他的主要成果都已经获得。然而，在 1870 年代，他的所有论文（除了最初的一篇简短笔记）都发表在挪威的学术期刊上，这阻碍了这些成果在欧洲其他地区的传播与认可。[4]1884 年，一位年轻的德国数学家弗里德里希·恩格尔来到李身边，与他一起撰写一部系统性著作，以全面阐述连续群理论。这一合作的成果就是三卷本的《变换群论》，分别于 1888 年、1890 年和 1893 年出版。“李群”这一术语最早于1893 年 出现，在李的学生阿图尔·特雷斯的博士论文中首次使用。[5]

李的思想并非孤立于数学整体之外。事实上，他对微分方程几何的兴趣最初是受到卡尔·古斯塔夫·雅可比工作的启发——雅可比曾研究一阶偏微分方程理论以及经典力学中的方程。雅可比的许多成果在他去世后的 1860 年代才得以发表，这在法国和德国引起了极大的兴趣。[6]李的执念（是建立一个微分方程对称性的理论，使其能完成与伽罗瓦在代数方程中所做的相同使命：即利用群论对它们进行分类。李和其他数学家证明，许多重要的特殊函数与正交多项式的方程，往往源于群论对称性。在李的早期工作中，他的目标是建立一个连续群的理论，以补充费利克斯·克莱因与昂利·庞加莱在模形式理论中所发展的离散群理论。李最初心中所想的应用领域，是微分方程理论。借鉴伽罗瓦理论及其在多项式方程中的作用，他的核心构想是：通过研究对称性，建立一个能够统一常微分方程整体领域的理论。然而，这一愿望并未完全实现。虽然对称性方法在常微分方程（ODE）中仍然是一个活跃的研究方向，但它并未成为该领域的主导工具。后来出现了微分伽罗瓦理论，但它是由其他数学家（如皮卡尔 Picard 和 维西奥 Vessiot）发展起来的，主要提供了一种求积理论，即通过不定积分来表达解的框架。

进一步推动人们思考连续群的动力，来自伯恩哈德·黎曼关于几何基础的思想，以及这些思想在克莱因（Klein）手中的进一步发展。由此，李在创建他的新理论时，结合了 19 世纪数学中的三大主题：

\begin{itemize}
\item 对称性的思想 —— 伽罗瓦通过群这一代数概念所展现的典范；
\item 几何理论与力学微分方程的显式解 —— 由泊松与雅可比所发展；
\item 几何学的新理解 —— 在普吕克、莫比乌斯、格拉斯曼等人的研究中逐步形成，并在黎曼的革命性愿景中达到顶峰。
\end{itemize}
虽然今天索弗斯·李被公认为连续群理论的奠基人，但在其结构理论的发展上，一个具有决定性意义的飞跃来自威廉·基林。他在1888 年发表了一系列题为《连续有限变换群的合成》 的论文中的第一篇，这一工作对后续数学的发展产生了深远影响。[7]

基林的成果后来被埃利·卡当加以精炼与推广，最终推动了半单李代数的分类、卡当的对称空间理论，以及赫尔曼·外尔利用最高权描述紧李群与半单李群的表示理论的建立。

1900年，大卫·希尔伯特在巴黎召开的国际数学家大会上提出了著名的希尔伯特第五星问题，向李群理论的研究者发出了挑战。

外尔（Hermann Weyl）推动了李群理论早期发展的成熟。他不仅完成了**半单李群不可约表示的分类**，并且将群论与量子力学建立了紧密联系；更重要的是，他通过清晰阐明**李的无穷小群（即李代数）**与**李群本身**之间的区别，使李的理论基础更加稳固，并开启了对李群**拓扑结构**的研究。[8]李群理论随后在现代数学语言的框架下被系统地重新表述，这一工作体现在克洛德·谢瓦莱所撰写的专著中。
\subsection{概述}